\documentclass[11pt,a4paper,twoside]{article}

% ── Encoding ──────────────────────────────────────────────────────────────────
\usepackage[utf8]{inputenc}
\usepackage[T1]{fontenc}
\usepackage[portuguese]{babel}

% ── Typography ────────────────────────────────────────────────────────────────
\usepackage{mathpazo}
\usepackage[scaled=0.90]{helvet}
\usepackage[activate={true,nocompatibility},final,tracking=true,kerning=true,expansion=false]{microtype}

% ── Layout ────────────────────────────────────────────────────────────────────
\usepackage[
  a4paper,
  top=3.0cm, bottom=3.2cm,
  inner=2.8cm, outer=2.4cm,
  headsep=1.0cm,
  footskip=1.2cm
]{geometry}

% ── Colors ────────────────────────────────────────────────────────────────────
\usepackage[table,dvipsnames]{xcolor}
\definecolor{INK}{HTML}{1A1A2E}
\definecolor{RULE}{HTML}{1856B8}
\definecolor{ACCENT}{HTML}{0C2A59}
\definecolor{MUTED}{HTML}{5C6475}
\definecolor{ROWSHADE}{HTML}{F4F6FA}
\definecolor{BOXBG}{HTML}{EEF2FA}
\definecolor{BOXLINE}{HTML}{1856B8}

% ── Tables ────────────────────────────────────────────────────────────────────
\usepackage{booktabs}
\usepackage{longtable}
\usepackage{tabularx}
\usepackage{array}
\usepackage{multirow}
\usepackage{makecell}
\renewcommand{\arraystretch}{1.34}
\setlength{\tabcolsep}{3.5pt}
\newcolumntype{Y}{>{\raggedright\arraybackslash}X}

% ── Lists ─────────────────────────────────────────────────────────────────────
\usepackage{enumitem}
\setlist{leftmargin=1.5em,itemsep=2pt,topsep=4pt,parsep=0pt}
\setlist[itemize,1]{label=\textcolor{RULE}{\rule[0.35ex]{4pt}{4pt}}}
\setlist[itemize,2]{label=\textcolor{MUTED}{\textendash}}
\setlist[enumerate,1]{label=\textcolor{RULE}{\arabic*.}}

% ── Section Titles ────────────────────────────────────────────────────────────
\usepackage{titlesec}
\titleformat{\section}
  {\sffamily\bfseries\large\color{RULE}}
  {\thesection}
  {0.8em}
  {}
  [\vspace{1pt}{\color{RULE}\titlerule[0.8pt]}\vspace{4pt}]

\titleformat{\subsection}
  {\sffamily\bfseries\normalsize\color{ACCENT}}
  {\thesubsection}
  {0.7em}
  {}

\titlespacing*{\section}{0pt}{20pt}{8pt}
\titlespacing*{\subsection}{0pt}{14pt}{4pt}

% ── Headers / Footers ─────────────────────────────────────────────────────────
\usepackage{fancyhdr}
\usepackage{lastpage}
\pagestyle{fancy}
\fancyhf{}
\fancyhead[LE,RO]{\small\sffamily\color{MUTED}Especificação Funcional - SMARTER HUB}
\fancyhead[RE,LO]{\small\sffamily\color{MUTED}PROES 2025/2026 - PC3}
\fancyfoot[C]{\small\sffamily\color{MUTED}\thepage\,/\,\pageref{LastPage}}
\renewcommand{\headrulewidth}{0.4pt}
\renewcommand{\headrule}{\color{RULE!40}\hrule height 0.4pt \vspace{0pt}}
\renewcommand{\footrulewidth}{0pt}

% ── Boxes / Callouts ──────────────────────────────────────────────────────────
\usepackage[most]{tcolorbox}
\tcbuselibrary{skins, breakable}

\newtcolorbox{callout}[1][] {
  enhanced, breakable,
  colback=BOXBG, colframe=BOXLINE,
  boxrule=0pt, leftrule=3pt,
  arc=0pt, outer arc=0pt,
  left=10pt, right=10pt, top=7pt, bottom=7pt,
  fontupper=\small,
  #1
}

\newcommand{\tH}[1]{\textbf{\sffamily\small\color{white}#1}}
\newcommand{\theadrow}{\rowcolor{ACCENT}}

% ── Cover image ──────────────────────────────────────────────────────────────
\usepackage{graphicx}
\usepackage{eso-pic}

% ── Misc ──────────────────────────────────────────────────────────────────────
\usepackage{setspace}
\usepackage{parskip}
\setlength{\parskip}{5pt}
\setlength{\parindent}{0pt}
\usepackage{hyperref}
\hypersetup{
  colorlinks=true,
  linkcolor=RULE,
  urlcolor=RULE,
  citecolor=RULE,
  pdftitle={Especificação Funcional SMARTER HUB - PC3},
  pdfauthor={Patrick Costa}
}

\begin{document}
\setstretch{1.17}

\begin{titlepage}
\thispagestyle{empty}
\AddToShipoutPictureBG*{%
  \includegraphics[width=\paperwidth,height=\paperheight]{capa_especificacao.png}%
}
\vspace*{\fill}
\end{titlepage}
\nopagecolor
\color{INK}

\tableofcontents
\newpage

\section{Objetivo, escopo e fronteiras}

\begin{callout}
\textbf{Objetivo do artefacto.} Consolidar, com nível de entrega académica e operacional, as regras funcionais do SMARTER HUB para o PC3, servindo como referência única para produto, UX, backend e testes. Este documento é de natureza evolutiva: deve ser atualizado sempre que uma alteração funcional for implementada, garantindo alinhamento contínuo entre especificação e produto.
\end{callout}

\subsection{Escopo funcional coberto}
\begin{itemize}
  \item Identidade e acesso: login local, Microsoft OAuth, sessão, perfil de acesso.
  \item Estrutura organizacional: equipas, subequipas, chefias, memberships, centros de custo.
  \item Ficha de membro e pedidos de alteração de dados; processo de admissão de novos membros.
  \item Férias e ausências: pedido, validação, aprovação, versionamento, calendário, exportação, ciclo processado/realizado.
  \item Banco de horas para membros BR, com saldo, relatório semanal PDF e política por estado (SP/RS).
  \item Plano de carreira individual: metas, competências e registo de evolução.
  \item Saúde e bem-estar: conteúdo informativo por perfil e notificações de medicina do trabalho.
  \item Permissões granulares, acesso total e trilho de auditoria.
  \item Formações, notificações, dashboard e assistente interno contextual.
  \item Internacionalização PT-PT / PT-BR com seleção de variante por utilizador.
\end{itemize}

\subsection{Fora de escopo nesta fase}
\begin{itemize}
  \item Decisões de implementação técnica de baixo nível (detalhe de código e infraestrutura).
  \item Otimizações avançadas sem impacto no comportamento funcional nucleado no PC2.
  \item Evoluções cosméticas de interface sem alteração de regra, fluxo ou autorização.
\end{itemize}

\subsection{Princípios mandatórios}
\begin{itemize}
  \item \textbf{Autorização antes de ação:} nenhuma operação crítica executa sem permissão efetiva e escopo válido.
  \item \textbf{Sem histórico perdido:} alterações a pedidos aprovados são realizadas por versionamento.
  \item \textbf{Sem rasto, sem commit:} decisão crítica exige auditoria persistida.
  \item \textbf{Regra por país:} políticas PT/BR são validadas antes de submissão e aprovação final.
  \item \textbf{Operação explicável:} bloqueios devem devolver causa e indicação de desbloqueio.
\end{itemize}

\section{Modelo de atores e autorização}

\subsection{Perfis funcionais}
\rowcolors{1}{ROWSHADE}{white}
\begin{longtable}{>{\small\bfseries}p{2.2cm} >{\small}p{4.1cm} >{\small}p{3.5cm} >{\small}p{3.8cm}}
\theadrow
\tH{Perfil} & \tH{Âmbito operacional} & \tH{Nível de decisão} & \tH{Limites estruturais} \\
\midrule
\endfirsthead
\theadrow
\tH{Perfil} & \tH{Âmbito operacional} & \tH{Nível de decisão} & \tH{Limites estruturais} \\
\midrule
\endhead
Membro & Operação individual (ficha, pedidos, notificações, formações) & Sem aprovação de terceiros por defeito & Atua no seu contexto e nas permissões que lhe forem atribuídas \\
Liderança & Gestão de equipa e aprovação de primeira linha & Aprova linha direta conforme cadeia ativa & Não atua fora do escopo de equipa autorizado \\
Liderança & Supervisão transversal e segunda linha de decisão & Aprova níveis superiores definidos na cadeia & Não substitui admin em governação global \\
Admin & Governação global de utilizadores, equipas e permissões & Pode atuar por exceção em PT quando aplicável & Deve manter rasto integral das ações críticas \\
Root Access & Superutilizador técnico-funcional & Último nível de autorização & Uso restrito a governação e manutenção sensível \\
\bottomrule
\end{longtable}

\subsection{Regras de aprovação e hierarquia}
\begin{enumerate}
  \item Um utilizador pode ter múltiplas memberships e papéis distintos por equipa.
  \item A aprovação base é efetuada pelos chefes das equipas relevantes do membro.
  \item Em múltiplas equipas, a aprovação é unânime entre aprovadores definidos para o pedido.
  \item Autoaprovação é sempre proibida.
  \item Rejeição de um aprovador encerra imediatamente o pedido como rejeitado.
  \item A lista de aprovadores é congelada no momento de submissão da versão.
  \item Sem equipa e sem acesso total, o fallback de aprovação é T.People.
  \item Com acesso total sem equipa, decide o concedente; indisponível, fallback T.People.
  \item Entre utilizadores com acesso total no mesmo nível hierárquico, não há decisão cruzada.
\end{enumerate}

\subsection{Matriz de capacidades por domínio}
\rowcolors{1}{ROWSHADE}{white}
\begin{longtable}{>{\small\bfseries}p{2.2cm} >{\small}p{4.1cm} >{\small}p{4.2cm} >{\small}p{3.2cm}}
\theadrow
\tH{Domínio} & \tH{Visualização} & \tH{Ação operacional} & \tH{Controlo/Auditoria} \\
\midrule
\endfirsthead
\theadrow
\tH{Domínio} & \tH{Visualização} & \tH{Ação operacional} & \tH{Controlo/Auditoria} \\
\midrule
\endhead
Identidade & Conta própria, perfil de acesso, sessão atual & Login, refresh de contexto, logout & Registo de autenticação e falhas \\
Ficha & Dados pessoais/contratuais conforme escopo & Submeter alterações e comprovativos & Histórico de pedido e motivo \\
Férias/Ausências & Histórico próprio e de equipa (quando permitido) & Submeter, editar pendente, cancelar, versionar aprovado & Trilha de decisões por etapa \\
Permissões & Permissões efetivas por utilizador autorizado & Grant/revoke, escopo, acesso total & Audit log de grants/revokes \\
Equipas & Estrutura e memberships por escopo & CRUD equipas (admin), gestão de membros & Rastreio de alterações estruturais \\
Formações & Plano próprio e organização (autorizados) & Criar, atribuir, concluir, exportar & Histórico de atribuição/conclusão \\
Dashboard & KPIs por equipa/organização e período & Filtrar cohort e exportar XLSX & Parâmetros de exportação registados \\
Notificações & Inbox pessoal e estados de leitura & Ler, limpar e navegar para ação & Integridade e idempotência \\
\bottomrule
\end{longtable}

\section{Regras de negócio de férias e ausências}

\subsection{Regras transversais}
\begin{itemize}
  \item Pedido inclui tipo, período, contexto de equipa e metadados de validação.
  \item Conflitos de período com pedidos \texttt{PENDING}/\texttt{APPROVED} bloqueiam submissão.
  \item Seleção com fim de semana bloqueia tipo \textit{Férias} e força \textit{Ausência}, quando aplicável.
  \item Pedido aprovado alterado gera nova versão pendente, sem sobrescrever versão anterior.
  \item Pedidos que atravessam o ano civil devem aparecer em ambos os anos relevantes.
  \item Meio-dia é suportado com indicação explícita manhã/tarde.
\end{itemize}

\subsection{Capacidade e calendário}
\begin{itemize}
  \item O sistema sinaliza risco quando a equipa se aproxima de ocupação crítica em férias simultâneas.
  \item A vista \textit{Férias da equipa} apresenta apenas pedidos aprovados.
  \item Distinção visual entre \textit{Férias} e \textit{Ausência} é obrigatória.
  \item Calendário inclui legenda de feriados, dias automáticos/extra e fins de semana.
\end{itemize}

\subsection{Catálogo de motivos de ausência}

O sistema oferece um catálogo padronizado de motivos para ausências justificadas, utilizado na classificação de pedidos e relatórios operacionais:

\begin{itemize}
  \item \textit{Não Justificada}
  \item \textit{Justificada - Greve}
  \item \textit{Justificada - Atividades empresa}
  \item \textit{Justificada - Assistência à família}
  \item \textit{Justificada - Estatuto trabalhador-estudante}
  \item \textit{Justificada - Licença casamento}
  \item \textit{Justificada - Consultas médicas}
  \item \textit{Justificada - Doença}
  \item \textit{Justificada - Formação}
  \item \textit{Justificada - Licença parental}
  \item \textit{Justificada - Morte de familiar}
  \item \textit{Justificada - Feriado local}
  \item \textit{Justificada - Obrigações legais}
  \item \textit{Justificada - Tolerância dada pela empresa}
\end{itemize}

\subsection{Políticas PT}
\begin{itemize}
  \item Aviso de conformidade para ausência de bloco de 10 dias úteis consecutivos (não bloqueante).
  \item Contratação no ano corrente: limite proporcional de 2 dias úteis/mês no 1.º ano (máximo 20).
  \item Bloqueio quando \texttt{PENDING + APPROVED + novo pedido} excede limite proporcional.
  \item Lembrete automático de transitados a expirar a 30 de abril (default: 30 dias antes).
\end{itemize}

\subsection{Políticas BR}
\begin{itemize}
  \item Início de férias não permitido à quinta-feira nem à sexta-feira.
  \item Início não permitido no dia útil imediatamente após feriado nacional.
  \item Validação do período concessivo com bloqueio por antecedência mínima inferior a 30 dias.
  \item Alteração/versionamento só com antecedência mínima de 10 dias sobre o início do pedido.
  \item Estagiário: bloqueado antes de 12 meses; após marco, proporcional de 30 dias/ano (2,5 dias/mês).
  \item Venda de férias (abono): até 1/3 do direito anual, teto absoluto de 10 dias.
  \item Aplicação diferenciada de feriados públicos por estado: São Paulo (SP) e Rio Grande do Sul (RS) utilizam calendários locais de feriados e pontos facultativos, incluindo apenas feriados globais ou feriados específicos do estado na validação de períodos de férias.
\end{itemize}

\subsection{Ciclo processado/realizado}
\begin{itemize}
  \item Após aprovação de férias, o sistema aguarda confirmação em dois momentos distintos.
  \item \textbf{Processado:} marcado pelo RH/aprovador para confirmar o processamento de pagamento das férias.
  \item \textbf{Realizado:} confirmado pelo membro após o regresso das férias, ou pelo aprovador em substituição.
  \item O sistema bloqueia nova confirmação de realizado para férias com data de fim futura.
  \item Ambas as confirmações registam o ator e o timestamp para audit log.
  \item O saldo definitivo só é atualizado após marcado como realizado.
\end{itemize}

\subsection{Banco de horas (BR)}
\begin{itemize}
  \item O banco de horas é aplicável a membros com \texttt{workCountry = BR}.
  \item Unidade oficial do saldo: \textbf{horas} (não minutos), com precisão de 2 casas decimais.
  \item Estrutura de saldo por membro: \textit{saldo creditado}, \textit{saldo debitado}, \textit{total} e \textit{limite}.
  \item Lançamentos manuais (crédito/débito) são operados por utilizadores com \textit{acesso total} (ou root), com motivo obrigatório.
  \item Limite acumulado operacional em BR: \textbf{100h} (controlo automático no cálculo e nos relatórios).
  \item Campo obrigatório de localização BR na ficha: \texttt{brWorkState} com opções \texttt{SP} e \texttt{RS}, utilizado para determinar calendário de feriados na validação de férias.
  \item Política de fecho por estado: \texttt{RS} semestral (abril/outubro); \texttt{SP} quadrimestral (fevereiro/junho/outubro).
  \item Excedente é calculado por \( |total| - limite \) quando positivo.
  \item Excedente gera aviso para o membro e para a chefia responsável, com cópia para utilizadores de acesso total.
  \item O sistema gera relatório semanal automático de banco de horas para acesso total, sem operação manual.
  \item O relatório semanal é gerado em \textbf{PDF} e armazenado em histórico próprio para consulta posterior (lista de semanas + abertura direta do ficheiro).
  \item Notificações e acesso à área de Banco de Horas são restritos a perfis com \texttt{workCountry = BR}; perfis PT não devem visualizar esta área.
  \item Chefias podem consultar visão por equipa; RH/acesso total podem consultar globalmente e exportar por equipa/geo.
\end{itemize}

\section{Experiência, vistas e ficheiro de membro}

\subsection{Navegação base por permissões}
\begin{itemize}
  \item \textbf{Home}: sempre visível.
  \item \textbf{A Minha Ficha}: autenticados (exceto caso técnico T.People).
  \item \textbf{Plano de Carreira}: autenticados (exceto T.People).
  \item \textbf{Dashboard}: apenas root/access total.
  \item \textbf{Aprovações}: requer pelo menos uma permissão de aprovação/visão global.
  \item \textbf{Admissões}: requer permissão de aprovação de perfil ou root/access total.
  \item \textbf{Férias}: requer permissão de pedido/consulta/gestão de regras.
  \item \textbf{Formações}: requer permissão de consulta/atribuição.
  \item \textbf{Banco de Horas}: apenas perfis BR (ou T.People) com permissão de consulta/gestão.
  \item \textbf{Saúde e Bem-estar}: não convidados.
\end{itemize}

\subsection{Seletor de variante linguística (PT-PT / PT-BR)}
\begin{itemize}
  \item O portal disponibiliza um seletor de variante linguística persistente no topo da interface, acessível em qualquer página autenticada.
  \item As opções são Portugal (PT-PT, por defeito) e Brasil (PT-BR), identificadas pela bandeira nacional respetiva.
  \item A preferência é guardada localmente (\texttt{localStorage}) e persiste entre sessões.
  \item A adaptação PT-BR abrange terminologia diferenciada (ex.: \textit{utilizador} \textrightarrow{} \textit{usuário}, \textit{guardar} \textrightarrow{} \textit{salvar}, \textit{equipa} \textrightarrow{} \textit{equipe}, \textit{formação} \textrightarrow{} \textit{treinamento}) e nomenclaturas de campos de ficha.
  \item Para conteúdo dinâmico (mensagens, descrições, erros), o sistema integra a API LibreTranslate como motor de tradução opcional, com cache local por texto.
\end{itemize}

\subsection{Vistas mandatórias para operação}
\begin{itemize}
  \item Separador \textit{Férias da equipa} com calendário mensal por membro e navegação temporal.
  \item Dashboard com filtros por período (incluindo presets plurianuais) e exportação XLSX.
  \item Utilizadores com acesso total podem exportar mapa de férias por intervalo livre.
\end{itemize}

\subsection{Ficha do membro e benefícios}
\begin{itemize}
  \item Secções: Dados Pessoais, Contacto, Documentos, Fiscal/Bancário, Emergência, Contratuais, Formação, Benefícios.
  \item Campo \textit{Número mecanográfico} obrigatório na área contratual.
  \item Centro de custo preenchido com base na equipa efetiva.
  \item Para \textit{workCountry = BR}: nomenclaturas locais (CPF, PIS, CTPS, RG, CNH, Título de eleitor, CEP), ocultação de campos não aplicáveis e sem bloco Cartão Continente.
  \item Voucher NOS: apenas contrato \textit{Sem termo}; cooldown de 2 anos entre pedidos.
\end{itemize}

\section{Assistente interno, qualidade e auditoria}

\subsection{Assistente interno (chatbot)}
\begin{itemize}
  \item Respostas contextualizadas por permissões efetivas e rota atual.
  \item Prioridade a instruções acionáveis: onde clicar, pré-condição em falta, motivo de bloqueio.
  \item Em falta de permissão, indica capacidade necessária para desbloqueio.
  \item Cobertura mínima: autenticação, ficha, benefícios, equipas, membros, permissões, aprovações, férias, formações, dashboard e notificações.
\end{itemize}

\subsection{Requisitos de qualidade e dados}
\begin{itemize}
  \item Integridade referencial entre utilizador, equipa, pedido, aprovação e permissão.
  \item Estados de pedido consistentes entre vistas e endpoints.
  \item Auditoria com autor, alvo, ação, timestamp e motivo (quando aplicável).
  \item Alteração funcional só considerada concluída com atualização desta especificação.
  \item Releases com evidência mínima de testes unitários, integração e E2E crítico.
\end{itemize}

\section{Fluxos ponta a ponta de referência}

\subsection{Fluxo A: pedido de férias}
\begin{enumerate}
  \item Membro submete pedido com período e contexto.
  \item Sistema valida regras transversais e política PT/BR.
  \item Sistema calcula aprovadores e congela cadeia.
  \item Aprovadores decidem por etapa.
  \item Pedido termina em aprovado, rejeitado, cancelado ou substituído por versão.
\end{enumerate}

\subsection{Fluxo B: alteração de ficha}
\begin{enumerate}
  \item Membro submete alteração de dados.
  \item Sistema cria pedido de alteração em estado pendente.
  \item Responsável aprova/rejeita com motivo.
  \item Sistema atualiza dados (se aprovado) e preserva histórico.
\end{enumerate}

\subsection{Fluxo C: governação de permissões}
\begin{enumerate}
  \item Gestor autorizado abre perfil de permissões do alvo.
  \item Concede/revoga permissão e escopo.
  \item Sistema regista evento de auditoria.
  \item Cache/contexto do alvo é reavaliado para refletir novo acesso.
\end{enumerate}

\subsection{Fluxo D: notificações operacionais}
\begin{enumerate}
  \item Evento de negócio gera notificação ao(s) utilizador(es) relevante(s).
  \item Utilizador consulta lista, marca como lida e/ou remove item.
  \item Em fluxos acionáveis, notificação navega para o ecrã de decisão.
\end{enumerate}

\subsection{Fluxo E: ciclo processado/realizado de férias}
\begin{enumerate}
  \item Férias transitam para estado \texttt{APPROVED}.
  \item RH/aprovador marca as férias como \textbf{processadas} (confirma pagamento/processamento).
  \item Após data de fim das férias, membro visualiza bann\-er de confirmação de realização.
  \item Membro confirma via modal, ou aprovador confirma em substituição.
  \item Sistema regista atores, timestamps e atualiza saldo definitivo.
\end{enumerate}

\subsection{Fluxo F: banco de horas (BR)}
\begin{enumerate}
  \item Membro BR visualiza saldo atual (creditado, debitado, total, limite).
  \item Utilizador com acesso total credita ou debita horas com motivo obrigatório.
  \item Sistema calcula excedente quando \( |total| > limite \) e envia aviso a membro, chefia e acesso total.
  \item Sistema gera relatório semanal automático em PDF, acessível em histórico.
  \item Chefias consultam visão por equipa; RH/acesso total exportam por estado/geo.
\end{enumerate}

\section{Rastreabilidade, riscos e critérios de aceite}

\subsection{Matriz de rastreabilidade}
\rowcolors{1}{ROWSHADE}{white}
\begin{longtable}{>{\small\bfseries}p{4.0cm} >{\small}p{4.7cm} >{\small}p{4.8cm}}
\theadrow
\tH{Tema pedido no PC2} & \tH{Cobertura nesta especificação} & \tH{Evidência de verificação} \\
\midrule
\endfirsthead
\theadrow
\tH{Tema pedido no PC2} & \tH{Cobertura nesta especificação} & \tH{Evidência de verificação} \\
\midrule
\endhead
Necessidade e problema & Escopo, princípios e regras de negócio & Fluxos definidos e políticas PT/BR explícitas \\
Contexto e solução & Modelo de atores, capacidades e navegação & Tabelas de domínio e autorização \\
Qualidade e governação & Requisitos de auditoria e dados & Critérios de aceite objetivos \\
Planeamento funcional & Definição de pronto e fronteiras de escopo & Itens mandatórios e risco monitorizado \\
\bottomrule
\end{longtable}

\subsection{Riscos funcionais monitorizados no PC2}
\rowcolors{1}{ROWSHADE}{white}
\begin{longtable}{>{\small\bfseries\centering\arraybackslash}p{0.9cm} >{\small}p{6.2cm} >{\small}p{6.0cm}}
\theadrow
\tH{ID} & \tH{Risco} & \tH{Mitigação definida} \\
\midrule
\endfirsthead
\theadrow
\tH{ID} & \tH{Risco} & \tH{Mitigação definida} \\
\midrule
\endhead
R1 & Alteração de regra sem atualização documental & Governação: alteração funcional só conclui com update desta especificação \\
R2 & Escalada indevida de privilégio em operação sensível & Validação por permissão efetiva, escopo e trilho de auditoria \\
R3 & Divergência entre hierarquia real e cadeia de aprovação & Revisão de memberships e cadeia congelada por versão \\
R4 & Regressão em regras PT/BR por alteração de lógica & Critérios de aceite por política e validação de conflitos \\
R5 & Perda de rastreabilidade de decisão & Obrigatoriedade de motivo, ator e timestamp em decisão crítica \\
\bottomrule
\end{longtable}

\subsection{Definição de pronto para entrega PC3}
\begin{itemize}
  \item Regras de aprovação e políticas PT/BR descritas sem contradições internas.
  \item Perfis, domínios e restrições alinhados com comportamento esperado do portal.
  \item Fluxos críticos mapeados com estados e pontos de auditoria.
  \item Vistas obrigatórias e restrições de acesso identificadas para UX e QA.
  \item Módulos Banco de Horas, Plano de Carreira, Admissões, Saúde e Bem-estar e Internacionalização cobertos.
  \item Documento mantido atualizado como referência oficial durante todo o período PC3.
\end{itemize}

\subsection{Critérios objetivos de aceitação deste artefacto}
\begin{itemize}
  \item \textbf{Completude}: domínios nucleares do PC3 cobertos sem lacunas estruturais.
  \item \textbf{Clareza}: regras críticas traduzíveis em comportamento observável de sistema.
  \item \textbf{Auditabilidade}: decisão de aprovação/revogação justificável por rasto de dados.
  \item \textbf{Testabilidade}: regras convertíveis em casos unitários, integração e E2E.
  \item \textbf{Evolução controlada}: qualquer alteração funcional é acompanhada de atualização deste documento.
\end{itemize}

\end{document}
